\documentclass{letter}
\usepackage{hyperref}
\date{\today}
\begin{document}

\begin{letter}{Editor-in-Chief \\
Journal of Statistical Mechanics: Theory and Experiment}

\opening{Dear Editor,}

I am pleased to submit my manuscript entitled ``Calibrating the Brody exponent
as a quantitative measure of short-range exclusion in 2D spatial point
processes'' for consideration for publication in the Journal of Statistical
Mechanics: Theory and Experiment.

\textbf{Why JSTAT.} The paper develops a calibrated statistical-mechanics
framework for short-range exclusion in 2D spatial point processes. The Brody
distribution---introduced as a phenomenological interpolation between Poisson
and Wigner level-spacing statistics in quantum chaos---is calibrated here
against 2D complete spatial randomness ($\beta=0.96\pm0.15$) and synthetic
hard-core and soft-core point processes. The resulting empirical
$\beta$--$r_{\text{excl}}$ calibration (Spearman $\rho=0.988$) translates
the dimensionless Brody exponent into a physically interpretable exclusion
measure. The framework is demonstrated on three statistically independent
classes of 2D point processes: 58 manufactured surface topographies,
phase-extracted interferometric profilometry, and binary embeddings of
prime numbers with sparse-integer and density-thinning controls. The paper
thus contributes directly to the statistical physics of spatial point
processes, random matrix theory applications, and exclusion models---all
core areas of JSTAT's scope.

\textbf{Key contributions.}
\begin{enumerate}
  \item \textbf{2D CSR baseline recalibration.} I demonstrate that the
  established 1D Poisson reference ($\beta=0$) is inappropriate for 2D
  point processes and recalibrate the complete-spatial-randomness baseline
  to $\beta=0.96\pm0.15$.

  \item \textbf{Empirical $\beta$--$r_{\text{excl}}$ calibration.} A monotonic
  relationship (Spearman $\rho=0.988$, $p<10^{-4}$) between the Brody exponent
  and the effective hard-core exclusion radius is established and validated,
  giving $\beta$ a direct physical interpretation.

  \item \textbf{Sparse-integer control for prime embeddings.} A novel control
  protocol proves that the Brody signal observed for 2D prime embeddings
  ($\beta=2.15$) is genuinely arithmetic rather than a trivial consequence
  of low point density ($\Delta\beta=+0.68$ over random-integer control).

  \item \textbf{Embedding-dependence demonstration.} A Cantor-pairing null
  result ($\beta=1.40$, TOST $p<0.01$) proves that 2D exclusion is created
  by the embedding geometry rather than being an intrinsic property of the
  prime sequence.

  \item \textbf{Experimental validation.} The framework is demonstrated on
  58 manufactured surfaces (10 materials, 10 manufacturing processes) and
  phase-extracted interferometric profilometry of a certified roundness
  standard ($\beta=2.00$ [1.90, 2.10]).
\end{enumerate}

\textbf{Significance.} The calibrated Brody framework provides the first
transferable, single-parameter measure of short-range exclusion strength
for 2D spatial point processes, with a validated physical interpretation
and a complete set of control protocols. The methods paper format ensures
that researchers across surface engineering, spatial statistics, and
mathematical physics can adopt the framework with confidence.

\textbf{Preprint and code availability.} A preprint of this manuscript is
available at \href{https://arxiv.org/abs/2606.16393}{arXiv:2606.16393}.
The complete analysis code and reproducibility materials are available at
\href{https://github.com/dawidkucharski/brody-2d-exclusion}{github.com/dawidkucharski/brody-2d-exclusion}.
Raw surface data (58 FV \texttt{.sur} files) and PSI interferogram frames are
available from the author upon request.

\textbf{Declaration.} This manuscript has not been published previously and
is not under consideration elsewhere. The author has approved the
manuscript and agrees with its submission to the Journal of Statistical Mechanics.
The author declares no competing interests.

\closing{Sincerely,}
Dawid Kucharski \\
Poznan University of Technology \\
Poznan, Poland \\
\href{mailto:dawid.kucharski@put.poznan.pl}{dawid.kucharski@put.poznan.pl}

\end{letter}
\end{document}
